% 17-agents.tex: AI agents, proof of personhood, and the future internet.
\section{AI agents, proof of personhood and the future internet}
\label{sec:agents}

This section was written as a statement of what did not exist. Most of it has since been built, and
it is rewritten here as a record of what runs, what it costs, and what is still missing, with the
marks distinguishing the three. The environment it describes has not changed; Polaris's answer to it
has.

\subsection{The problem that is arriving}

When software agents can generate unlimited accounts, content, transactions, messages and apparent
participants, the cheap signal that a request came from a person stops being available for free.
Many digital systems will need to distinguish, without learning more than they must, between one
accountable human; one authorized agent acting for that human or for an institution; anonymous
machine activity; duplicated or synthetic identities; and organisational agents acting under an
institution's authority. Every naive way of obtaining that signal leaks identity: accounts, phone
numbers, payment instruments, biometrics. The identity question changes shape. It is no longer only
\emph{who are you}, but \emph{what kind of actor are you, whose authority are you exercising, what
are you allowed to do, and how can I verify all of that without learning anything else}. That is
the future-facing question this paper considers the most important one Polaris's core could be
extended to answer.

\subsection{What exists today that points this way}

Six ingredients of a privacy-preserving proof of personhood already run; the last three arrived
after this section was first written.

\begin{itemize}
  \item \textbf{Uniqueness at issuance.} C3, one active token per person, is a partial unique index
    in the database. That is the root of Sybil resistance: a human can hold one active credential, and
    no application code can bypass the constraint. \sImpl
  \item \textbf{Membership without identity.} The epoch membership proof proves that a valid token
    exists in an authority's committed set with public inputs of only the root, the epoch, a context
    and a nonce; a zero-knowledge verification records no token at all; the anonymity set is a full
    epoch. \sImpl
  \item \textbf{Revocation without naming.} A revoked token drops out of the next epoch root, so a
    membership proof expires with the epoch and no relying party learns who was revoked. \sImpl
  \item \textbf{A key the holder controls.} A credential may carry a holder public key the issuer
    binds to it, so possession of the file stops being possession of the credential and a verifier
    can require a signature from the person rather than a copy of their credential. \sImpl
  \item \textbf{Proving on the holder's device.} The epoch's leaf set is published, signed and
    bounded; every requester receives identical bytes; the holder finds their own leaf and builds the
    path locally. The issuer is never told which member is proving. \sImpl
  \item \textbf{A per-verifier nullifier.} The proof carries
    \code{Poseidon(secret\,||\,scope\,||\,epoch)}, constant for one person within one relying
    party's scope and epoch and uncorrelated across scopes. \sImpl
\end{itemize}

\subsection{One human, once: what runs and what it costs}

The membership proof stops replay of one proof, but on its own it does not stop the same person
producing a thousand fresh proofs to one site and registering a thousand accounts. \emph{One human,
once, here} needs a scoped nullifier: one more public input, a hash of the holder's secret, the
relying party's scope and the epoch, so a site can refuse a second proof from the same person while
two sites cannot link their nullifiers. That is now a public input of the circuit. \sImpl

Adding it was not the small change this section once predicted, and the reason is worth recording
because it is the kind of mistake a plan makes from outside the code. The epoch leaf was a
\code{SHA3-256} digest computed \emph{outside} the circuit and handed in as an opaque private
value. A nullifier beside an opaque leaf proves only that the prover knows some number: nothing ties
the two to one secret, so a prover could pair any member's leaf with a nullifier of their own
choosing, and every property above would have been a claim rather than a proof. Constraining both to
one secret means opening the leaf inside the circuit, and verifying a \code{SHA3-256} preimage there
costs thousands of constraints where Poseidon costs about a hundred. So the leaf became
\code{Poseidon(secret\,||\,context)}, a commitment the circuit opens, and the old digest became the
holder's secret. The change is not backward compatible: an epoch closed under the old derivation
holds leaves the current circuit cannot open, and epochs are re-closed rather than migrated.

The two conditions the vocation imposes were carried into the implementation rather than left as
intentions. The nullifier's scope is per relying party and per epoch, never global, because a global
nullifier is a persistent identifier and therefore the surveillance outcome the constitution refuses;
and it rotates each epoch, so membership never becomes a permanent pseudonym. A relying party's
ledger is keyed by its own scope and that epoch, and should not outlive it.

Two bounds remain and are stated rather than closed. The nullifier does not hide the holder from the
\emph{issuer}, which derives every leaf in order to build the tree and could therefore compute any
member's nullifier in any scope; the property holds between relying parties. And the presentation
layer is still full-credential: a verifier looking at a plain credential sees a stable token value,
an issuer signature and a holder key, so two verifiers who deliberately keep the raw material can
still correlate. What the pairwise derivation bounds is what a verifier \emph{stores}, which is the
material that gets pooled, sold, subpoenaed and breached. Section~\ref{sec:privacy} states the
distinction, and the verifier reports which of the two a given presentation gave rather than letting
a caller assume the stronger one. The blinded and one-time presentation technologies that would
remove the residue are surveyed in Section~\ref{sec:research}; none is implemented. \sFut

\subsection{Delegating authority to an agent}

\figfile{delegation}{The delegation chain, as built. The human holds the foundational credential and
a key the issuer bound to it. With that key they sign a grant naming actions, limits, an expiry and a
revocation handle; the agent signs each action against the service's own nonce; the service verifies
the chain offline. The one link drawn dashed is the residue: under a plain holder binding the service
also sees the credential's token value, so the correlation it is left with is bounded rather than
absent.}{fig:delegation}

The relationship is a chain, and it now runs. A human possesses a foundational credential and a key
the issuer bound to it. With that key the human signs a grant giving an agent a narrowly scoped
authority: named actions, stated limits, an expiry, and a revocation handle belonging to the grant
alone. The agent signs each action it performs, naming the action and the service's own nonce. A
service verifies, offline, that the grant is genuine and unexpired, that it covers this exact action,
that it has not been revoked, and that it chains to a holder key an issuer bound to a real
credential. \sImpl

What the design is for is visible in what it refuses. Handing an agent the credential itself gives it
everything the person can do, forever, revocable only by revoking the person. The grant is the
opposite of each: \code{actions} and \code{limits} sit \emph{inside} the signed statement, so a
grant widened in transit fails rather than passing invisibly; an empty action list authorises nothing
and there is no way to write \emph{everything}; a limit key the service does not recognise is
refused rather than ignored, since a grant that says \code{max\_transfers: 3} to a service that has
never heard of it must not read as unlimited. And the agent must sign: without that link a grant is a
bearer token, and whoever copies it in transit becomes the agent.

Revocation is the property worth naming twice. The revocation is signed by the same key that signed
the grant, so anyone may publish bytes but only the holder may end it. The issuer is not contacted
and never learns the grant existed, and the human's credential is untouched throughout. A person can
end their agent's authority without asking permission from, or being observed by, the authority that
issued their identity.

The revocation carries the grant identifier and the instant, and nothing else. It will not carry a
reason. A field recording \emph{why} a grant ended is a field a coercer can demand be filled in or
left empty, and either way it turns a revocation into a signal about the person; the constitutional
argument is the same one that keeps the presented code out of the holder proof's signed statement.

What a service learns is bounded, and the bound is reported rather than assumed. Under a plain holder
binding the service also sees the credential's token value, so the verdict names the correlation
\code{exposed}; only a chain carried by the zero-knowledge path, whose handle is the scoped
nullifier, is named \code{bounded}. The grant hides the human from the agent's \emph{actions}. It
does not hide the credential from the service.

\subsection{The future internet, concretely}

The term Web3 is used here in one narrow sense, and it has nothing to do with currency: an internet
in which people, autonomous agents, institutions, wallets, distributed services and independently
operated networks need portable cryptographic trust that does not route through one company or one
state. Against that environment the following table marks what Polaris could provide and what each
item's status is today.

{\footnotesize
\setlength{\tabcolsep}{4pt}\renewcommand{\arraystretch}{1.12}
\rowcolors{2}{tabstripe}{white}
\begin{xltabular}{\textwidth}{>{\RaggedRight\arraybackslash}p{1.85in} >{\RaggedRight\arraybackslash}X >{\RaggedRight\arraybackslash}p{1.8in}}
\caption{What a future environment would ask of Polaris, and the status of each.}\label{tab:web3}\\
\toprule
\rowcolor{tabheader}\thead{The environment asks for} & \thead{What Polaris has} & \thead{Status} \\
\midrule
\endfirsthead
\toprule
\rowcolor{tabheader}\thead{The environment asks for} & \thead{What Polaris has} & \thead{Status} \\
\midrule
\endhead
\midrule\multicolumn{3}{r}{\itshape continued on the next page}\\
\endfoot
\bottomrule
\endlastfoot
Portable credentials & A signed file any standard library verifies & \sImpl \\
Independently verifiable proofs & The detached verifier, two SDKs, a conformance suite & \sImpl\newline\sExt \\
Institutional federation & Per-authority roots, directional attestations, published feeds & \sImpl\newline\sExt \\
Service discovery & The signed registry & \sImpl \\
Document signatures with long-term validity & Signed containers with evidence fixed at the instant & \sImpl \\
Transparent authority behaviour & Three logs, witnesses, equivocation proofs & \sImpl\newline\sExt \\
Offline verification & Status assertions stapled to presentations & \sImpl \\
Post-quantum trust & ML-DSA, hash-based proofs, a hybrid edge; classical remainders named & \sImpl \\
Human, account and agent authorization relationships & A login token whose subject is per relying party, plus the agent grant below & \sImpl \\
Privacy-preserving uniqueness & Membership proofs carrying a per-verifier, per-epoch nullifier & \sImpl \\
Delegation to agents with bounded scope & A holder-signed grant: actions, limits, expiry, revocation handle & \sImpl \\
Revoking an agent without destroying the human's identity & A holder-signed revocation; the issuer is not contacted and the credential is untouched & \sImpl \\
Presentation that shows a verifier nothing stable & The zero-knowledge path only; a plain credential still shows a stable value & \sImpl\newline\sFut \\
\end{xltabular}
}
